% =============================================================
%  CS102-F26 Digital Design and Computer Architecture I
%  Lecture 13 -- Counters (Unit 13)
%  Follows Section 5.4 (5.4.1 and 5.4.2) of
%  A Student's Computer Architecture Book I.
%
%  Build the slides and the print handout:  npm run materials -- lecture-13
% =============================================================
\makeatletter
\def\input@path{{templates/latex/}{../}}
\makeatother

\documentclass[10pt]{wosm-slides}
\usepackage{cs102-lectures}[2026/10/20]

\title[Lecture 13 \textperiodcentered\ Counters]{Counters}
\subtitle{Lecture 13: the synchronous up counter and the asynchronous ripple counter}

% ---------- Helpers for this lecture --------------------------
% \pic (f) at (<x>,<y>) {tff};  T flip-flop with a reset, as in the book:
% (f-t), (f-clk), (f-q), and (f-r) at the bottom.
\tikzset{
  tff/.pic={
    \draw[wosm-text, line width=0.6pt] (-0.5,-0.55) rectangle (0.5,0.55);
    \draw[wosm-text, line width=0.6pt] (-0.5,-0.15) -- (-0.36,-0.27) -- (-0.5,-0.39);
    \node[font=\scriptsize, anchor=west] at (-0.5,0.27) {T};
    \node[font=\scriptsize, anchor=east] at (0.5,0.27) {Q};
    \node[font=\tiny, anchor=south] at (0,-0.55) {R};
    \coordinate (-t) at (-0.5,0.27);
    \coordinate (-clk) at (-0.5,-0.27);
    \coordinate (-q) at (0.5,0.27);
    \coordinate (-r) at (0,-0.55);
  },
  fsm/.style={st, minimum size=0.72cm, inner sep=1pt, font=\scriptsize},
}
\newcommand*\junction[1]{\fill[wosm-text] #1 circle (1.5pt);}

% =============================================================
\begin{document}

\maketitle

\begin{frame}{Today}
  \tableofcontents
\end{frame}

% =============================================================
\section{Counting}

\begin{frame}{What a counter is}
  \begin{columns}[T, totalwidth=\textwidth]
    \column{0.5\textwidth}
      \begin{definition}[Counter]
        An $n$-bit \emph{counter} outputs an $n$-bit binary number and, at each
        clock edge, adds 1 to it (an \emph{up counter}) or takes 1 away (a
        \emph{down counter}).
      \end{definition}
      \begin{itemize}
        \item 0\ldots\ 1\ldots\ 2\ldots\ 3\ldots: the state is the number
          shown so far.
        \item Between rising edges the output holds, as in every flip-flop.
      \end{itemize}
    \column{0.44\textwidth}
      \centering
      \small
      \begin{tabular}{@{}c|c|cc@{}}
        clk & $Q_t$ & $Q^{\mathrm{up}}_{t+1}$ & $Q^{\mathrm{down}}_{t+1}$ \\
        \midrule
        $\uparrow$ & 000 & 001 & 111 \\
        $\uparrow$ & 001 & 010 & 000 \\
        $\uparrow$ & 010 & 011 & 001 \\
        $\uparrow$ & 011 & 100 & 010 \\
        $\uparrow$ & 100 & 101 & 011 \\
        $\uparrow$ & 101 & 110 & 100 \\
        $\uparrow$ & 110 & 111 & 101 \\
        $\uparrow$ & 111 & 000 & 110 \\
        \midrule
        0, 1, $\downarrow$ & $q$ & $q$ & $q$ \\
      \end{tabular}\par\smallskip
      \wosmlabel{A 3-bit up and down counter (Table 5.15)}
  \end{columns}
\end{frame}

\begin{frame}{Wrapping around}
  \begin{columns}[c, totalwidth=\textwidth]
    \column{0.6\textwidth}
      \begin{itemize}
        \item After 111 (7), the up counter goes back to 000. After 000, the down
          counter goes to 111.
        \item So a 3-bit up counter that starts at 0 shows $k \bmod 8$ after
          $k$ edges; an $n$-bit one counts modulo $2^n$.
        \item As a state machine \unitref{10}: 8 states in a ring, one step per
          clock edge, no other input. The output is the state.
      \end{itemize}
    \column{0.36\textwidth}
      \centering
      \begin{tikzpicture}
        \foreach \i/\b in {0/000, 1/001, 2/010, 3/011, 4/100, 5/101, 6/110, 7/111}
          {\node[fsm] (s\i) at ({90-45*\i}:1.45) {\b};}
        \foreach \i [evaluate=\i as \j using {int(mod(\i+1,8))}] in {0,...,7}
          {\draw[arrow] (s\i) -- (s\j);}
      \end{tikzpicture}
  \end{columns}

  \begin{remark}[Not a bug]
    The wrap-around is a feature. The DX8's stepper \unitref{21} is a 3-bit counter
    with a decoder: it goes through steps 1 to 8, then starts again at step 1.
  \end{remark}
\end{frame}

% =============================================================
\section{Synchronous up counter}

\begin{frame}{Which bits flip?}
  \begin{columns}[c, totalwidth=\textwidth]
    \column{0.4\textwidth}
      \centering
      \begin{tabular}{@{}ccc|ccc@{}}
        $y_2$ & $y_1$ & $y_0$ & $T_2$ & $T_1$ & $T_0$ \\
        \midrule
        0 & 0 & 0 & 0 & 0 & 1 \\
        0 & 0 & 1 & 0 & 1 & 1 \\
        0 & 1 & 0 & 0 & 0 & 1 \\
        0 & 1 & 1 & 1 & 1 & 1 \\
        1 & 0 & 0 & 0 & 0 & 1 \\
        1 & 0 & 1 & 0 & 1 & 1 \\
        1 & 1 & 0 & 0 & 0 & 1 \\
        1 & 1 & 1 & 1 & 1 & 1 \\
      \end{tabular}\par\smallskip
      \wosmlabel{$T_i = 1$: bit $i$ flips at the next edge}
    \column{0.54\textwidth}
      \begin{itemize}
        \item Adding 1 in binary: bit $i$ flips exactly when all the bits below it
          are 1 (the carry reaches it).
        \item A T flip-flop flips when $T = 1$ \unitref{12}. So use one per bit,
          with
          \[ T_0 = 1, \qquad T_1 = y_0, \qquad T_2 = y_0\, y_1. \]
        \item In general, $T_i = y_0\, y_1 \cdots y_{i-1}$.
      \end{itemize}
  \end{columns}
\end{frame}

\begin{frame}{The circuit}
  \begin{center}
    \begin{tikzpicture}[gates]
      \foreach \i/\x in {0/0, 1/2.8, 2/5.8} {\pic (f\i) at (\x,0) {tff};}
      \foreach \i in {0,1,2} {
        \draw[wire] (f\i-q) -- ++(0.3,0) coordinate (y\i) -- ++(0,1.25)
          node[above] {$y_\i$};
      }
      % T inputs
      \draw[wire] (f0-t) -- ++(-0.4,0) node[left] {1};
      \draw[wire] ($(y0)+(0,0.7)$) -| ($(f1-t)+(-0.3,0)$) -- (f1-t);
      \node[and gate, scale=0.8] (a) at (4.6,0.6) {};
      \draw[wire] ($(f1-t)+(-0.3,0.7)$) -| ($(a.input 1)+(-0.2,0)$) -- (a.input 1);
      \draw[wire] (y1 |- a.input 2) -- (a.input 2);
      \draw[wire] (a.output) -- ++(0.12,0) |- (f2-t);
      \junction{($(y0)+(0,0.7)$)}
      \junction{($(f1-t)+(-0.3,0.7)$)}
      \junction{(y1 |- a.input 2)}
      % clock
      \draw[wire] (-1.2,-0.95) node[left] {clk} -| ($(f2-clk)+(-0.25,0)$) -- (f2-clk);
      \foreach \i in {0,1} {
        \draw[wire] ($(f\i-clk)+(-0.25,-0.68)$) |- (f\i-clk);
        \junction{($(f\i-clk)+(-0.25,-0.68)$)}
      }
      % reset
      \draw[wire] (-1.2,-1.3) node[left] {$R$} -| (f2-r);
      \foreach \i in {0,1} {
        \draw[wire] (f\i-r) -- (f\i-r |- 0,-1.3);
        \junction{(f\i-r |- 0,-1.3)}
      }
    \end{tikzpicture}
  \end{center}

  \begin{itemize}
    \item One clock for all three flip-flops: they all change at the same edges,
      \alert{in sync} with the clock. Hence \emph{synchronous} counter.
    \item $R$ is the asynchronous reset of Lecture 12: $R = 1$ sets the count to 000.
  \end{itemize}
\end{frame}

\begin{frame}{In time}
  \begin{center}
    \begin{tikzpicture}[x=0.55cm]
      \foreach \t in {1,3,...,15} {\timingedge{\t}{3}}
      \timingwave{0}{clk}{0,1,0,1,0,1,0,1,0,1,0,1,0,1,0,1}
      \timingwave{1}{$y_0$}{0,1,1,0,0,1,1,0,0,1,1,0,0,1,1,0}
      \timingwave{2}{$y_1$}{0,0,0,1,1,1,1,0,0,0,0,1,1,1,1,0}
      \timingwave{3}{$y_2$}{0,0,0,0,0,0,0,1,1,1,1,1,1,1,1,0}
      \node[left, font=\scriptsize] at (-0.1,-1.94) {count};
      \foreach \x/\v in {0.5/0, 2/1, 4/2, 6/3, 8/4, 10/5, 12/6, 14/7, 15.5/0}
        {\node[font=\scriptsize, text=wosm-brick] at (\x,-1.94) {\v};}
    \end{tikzpicture}
  \end{center}

  \begin{itemize}
    \item All the bits that flip, flip together, right after a rising edge.
    \item $y_0$ has half the frequency of the clock, $y_1$ a quarter, $y_2$ an eighth:
      bit $i$ divides the clock frequency by $2^{i+1}$.
  \end{itemize}
\end{frame}

\begin{frame}{More bits}
  \begin{itemize}
    \item A 4-bit counter adds $T_3 = y_0\, y_1\, y_2$; an 8-bit one goes on to
      $T_7 = y_0\, y_1 \cdots y_6$.
    \item Each $T_i$ is the one before it, ANDed with one more bit:
      $T_{i} = T_{i-1}\, y_{i-1}$. A chain of AND gates.
    \item With D flip-flops instead, feed each one $D_i = y_i \oplus T_i$: that is
      $Q^{+} = T \oplus Q$ again \unitref{12}.
  \end{itemize}

  \begin{remark}[The longest path]
    After an edge, the new $y_0$ has to pass through the whole AND chain before
    $T_{n-1}$ is ready for the next edge. Lecture 14 turns this into a limit on the
    clock period.
  \end{remark}
\end{frame}

% =============================================================
\section{Asynchronous ripple counter}

\begin{frame}{Clocked by the bit before}
  \begin{center}
    \begin{tikzpicture}[gates]
      \foreach \i/\x in {0/0, 1/2.8, 2/5.6} {\pic (f\i) at (\x,0) {tff};}
      \foreach \i in {0,1,2} {
        \draw[wire] (f\i-t) -- ++(-0.35,0) node[left, font=\scriptsize] {1};
        \draw[wire] (f\i-q) -- ++(0.3,0) coordinate (y\i) -- ++(0,1.0)
          node[above] {$y_\i$};
      }
      \foreach \i/\j in {0/1, 1/2} {
        \node[not gate, scale=0.7] (n\j) at ($(f\j-clk)+(-0.85,0)$) {};
        \draw[wire] (y\i) |- (n\j.input);
        \draw[wire] (n\j.output) -- (f\j-clk);
        \junction{(y\i)}
      }
      \draw[wire] (-1.2,-0.95) node[left] {clk} -| ($(f0-clk)+(-0.25,0)$) -- (f0-clk);
      \draw[wire] (-1.2,-1.3) node[left] {$R$} -| (f2-r);
      \foreach \i in {0,1} {
        \draw[wire] (f\i-r) -- (f\i-r |- 0,-1.3);
        \junction{(f\i-r |- 0,-1.3)}
      }
    \end{tikzpicture}
  \end{center}

  \begin{itemize}
    \item Every $T$ is 1: each flip-flop flips at each rising edge of \emph{its own}
      clock. Only the first one gets the real clock; flip-flop $i$ gets
      $\overline{y_{i-1}}$.
    \item $\overline{y_{i-1}}$ rises exactly when $y_{i-1}$ falls from 1 to 0: when
      a carry moves up into bit $i$.
  \end{itemize}
\end{frame}

\begin{frame}{Step by step}
  \begin{center}
    \begin{tabular}{@{}l|ccccccccc@{}}
              & $t = 0$ & 1 & 2 & 3 & 4 & 5 & 6 & 7 & 8 \\
      \midrule
      $\mathrm{clk}_0 = \mathrm{clk}$
              &  & $\uparrow$ & $\uparrow$ & $\uparrow$ & $\uparrow$ & $\uparrow$
              & $\uparrow$ & $\uparrow$ & $\uparrow$ \\
      $y_0$   & 0 & 1 & 0 & 1 & 0 & 1 & 0 & 1 & 0 \\
      $\mathrm{clk}_1 = \overline{y_0}$
              & 1 & 0 & $\uparrow$ & 0 & $\uparrow$ & 0 & $\uparrow$ & 0 & $\uparrow$ \\
      $y_1$   & 0 & 0 & 1 & 1 & 0 & 0 & 1 & 1 & 0 \\
      $\mathrm{clk}_2 = \overline{y_1}$
              & 1 & 1 & 0 & 0 & $\uparrow$ & 1 & 0 & 0 & $\uparrow$ \\
      $y_2$   & 0 & 0 & 0 & 0 & 1 & 1 & 1 & 1 & 0 \\
      \midrule
      $y_2 y_1 y_0$ & 0 & 1 & 2 & 3 & 4 & 5 & 6 & 7 & 0 \\
    \end{tabular}\par\smallskip
    \wosmlabel{Analysis of the 3-bit ripple counter (Table 5.16); $\uparrow$: a rising
    edge}
  \end{center}

  \begin{itemize}
    \item $y_0$ flips at every clock edge, $y_1$ at every second one, $y_2$ at every
      fourth.
    \item Each stage clocks the next one: the clock edge \emph{ripples} through the
      flip-flops. Hence \emph{ripple counter}; and since they do not change at the
      same time, \emph{asynchronous}.
  \end{itemize}
\end{frame}

\begin{frame}{The ripple}
  \begin{center}
    \begin{tikzpicture}[x=0.135cm]
      \foreach \t in {4,12,...,60} {\timingedge{\t}{3}}
      \timingwave{0}{clk}{0,0,0,0,1,1,1,1,0,0,0,0,1,1,1,1,0,0,0,0,1,1,1,1,0,0,0,0,1,1,1,1,0,0,0,0,1,1,1,1,0,0,0,0,1,1,1,1,0,0,0,0,1,1,1,1,0,0,0,0,1,1,1,1,0,0}
      \timingwave{1}{$y_0$}{0,0,0,0,0,1,1,1,1,1,1,1,1,0,0,0,0,0,0,0,0,1,1,1,1,1,1,1,1,0,0,0,0,0,0,0,0,1,1,1,1,1,1,1,1,0,0,0,0,0,0,0,0,1,1,1,1,1,1,1,1,0,0,0,0,0}
      \timingwave{2}{$y_1$}{0,0,0,0,0,0,0,0,0,0,0,0,0,0,1,1,1,1,1,1,1,1,1,1,1,1,1,1,1,1,0,0,0,0,0,0,0,0,0,0,0,0,0,0,0,0,1,1,1,1,1,1,1,1,1,1,1,1,1,1,1,1,0,0,0,0}
      \timingwave{3}{$y_2$}{0,0,0,0,0,0,0,0,0,0,0,0,0,0,0,0,0,0,0,0,0,0,0,0,0,0,0,0,0,0,0,1,1,1,1,1,1,1,1,1,1,1,1,1,1,1,1,1,1,1,1,1,1,1,1,1,1,1,1,1,1,1,1,0,0,0}
      \node[left, font=\scriptsize] at (-0.1,-1.94) {count};
      \foreach \x/\v in {2.5/0, 9/1, 17.5/2, 25/3, 34/4, 41/5, 49.5/6, 57/7, 64.5/0}
        {\node[font=\scriptsize, text=wosm-brick] at (\x,-1.94) {\v};}
    \end{tikzpicture}
  \end{center}

  \begin{itemize}
    \item Real flip-flops take a little time to change their output (drawn here much
      larger than it is). In a ripple counter these delays add up: $y_1$ changes one
      delay after $y_0$, $y_2$ one delay after $y_1$.
    \item With $n$ flip-flops, the count is correct only $n$ delays after the edge.
  \end{itemize}
\end{frame}

\begin{frame}{From 3 to 4, slowly}
  \begin{columns}[c, totalwidth=\textwidth]
    \column{0.5\textwidth}
      \centering
      \begin{tikzpicture}[x=0.95cm]
        \foreach \t in {1} {\timingedge{\t}{3}}
        \timingwave{0}{clk}{0,1,1,1,1,0}
        \timingwave{1}{$y_0$}{1,1,0,0,0,0}
        \timingwave{2}{$y_1$}{1,1,1,0,0,0}
        \timingwave{3}{$y_2$}{0,0,0,0,1,1}
        \node[left, font=\scriptsize] at (-0.1,-1.94) {count};
        \foreach \x/\v in {1/3, 2.5/2, 3.5/0, 5/4}
          {\node[font=\scriptsize, text=wosm-brick] at (\x,-1.94) {\v};}
        \foreach \t in {2,3,4} {\draw[wosm-rule, dashed, line width=0.4pt]
          (\t,0.4) -- (\t,-1.8);}
      \end{tikzpicture}
    \column{0.46\textwidth}
      \begin{itemize}
        \item 011 $\to$ 010: $y_0$ falls.
        \item 010 $\to$ 000: that clocks $y_1$, which falls.
        \item 000 $\to$ 100: that clocks $y_2$, which rises.
        \item For a moment the counter shows 2, then 0, before 4. A circuit that
          reads it too early gets a wrong number, like the glitches of
          \unitref{9}.
      \end{itemize}
  \end{columns}
\end{frame}

\begin{frame}{Synchronous or ripple?}
  \begin{center}
    \begin{tabular}{@{}lll@{}}
      \toprule
       & Synchronous & Ripple \\
      \midrule
      Clock of flip-flop $i$ & clk, for all & clk for the first, then
        $\overline{y_{i-1}}$ \\
      $T$ inputs & $T_i = y_0 \cdots y_{i-1}$ & all 1 \\
      Extra gates & a chain of AND gates & NOT gates only \\
      Outputs change & together & one after another \\
      Correct after an edge & one flip-flop delay & $n$ delays \\
      Wrong values on the way & no & yes \\
      \bottomrule
    \end{tabular}
  \end{center}

  \begin{remark}
    The ripple counter is the smaller circuit; the synchronous counter is the one
    the rest of a synchronous design can trust at every clock edge. The DX8's
    stepper uses the synchronous one.
  \end{remark}
\end{frame}

\begin{frame}{Try it}
  \begin{enumerate}
    \item A 3-bit up counter starts at 000. What does it show after 13 rising
      edges?\par
      \uncover<2->{\small 101: $13 \bmod 8 = 5$.}
    \item Which $T$ inputs make a 4-bit synchronous up counter?\par
      \uncover<3->{\small $T_0 = 1$, $T_1 = y_0$, $T_2 = y_0 y_1$,
        $T_3 = y_0 y_1 y_2$.}
    \item Make the synchronous counter count down.\par
      \uncover<4->{\small Bit $i$ flips when all the bits below it are 0:
        $T_0 = 1$, $T_1 = \overline{y_0}$, $T_2 = \overline{y_0}\,\overline{y_1}$.}
    \item Make the ripple counter count down.\par
      \uncover<5->{\small Clock flip-flop $i$ with $y_{i-1}$ instead of
        $\overline{y_{i-1}}$: 000, 111, 110, \ldots}
    \item In a 4-bit ripple counter, each flip-flop and its NOT gate take 10\,ns.
      What happens on the step from 0111 to 1000?\par
      \uncover<6->{\small 0111, 0110, 0100, 0000, 1000: correct after 40\,ns.}
  \end{enumerate}
\end{frame}

% =============================================================
\section{Summary}

\begin{frame}{Summary}
  \begin{itemize}
    \item An $n$-bit counter adds (or takes away) 1 at each clock edge and wraps
      around: it counts modulo $2^n$.
    \item When counting up, bit $i$ flips exactly when all the bits below it are 1.
    \item Synchronous counter: one clock for all, $T_i = y_0 \cdots y_{i-1}$; all
      bits change together.
    \item Ripple counter: every $T = 1$, each flip-flop clocked by
      $\overline{y_{i-1}}$; fewer gates, but the delays add up and wrong values
      appear on the way.
    \item Bit $i$ of a counter divides the clock frequency by $2^{i+1}$.
  \end{itemize}
\end{frame}

\begin{frame}{Before the next lecture}
  \begin{itemize}
    \item Read Section~5.4 of the book.
    \item In Logisim-evolution, build both 3-bit counters from your T flip-flops
      of Lecture 12, with a reset, and check that they count from 0 to 7.
    \item Put a 3-to-8 decoder \unitref{8} after the synchronous counter: one output
      at a time is 1, in turn. That is the DX8's stepper \unitref{21}.
  \end{itemize}

  \begin{block}{Next lecture \unitref{14}}
    Timing of sequential circuits: the setup and hold time constraints.
  \end{block}
\end{frame}

\end{document}
